\chapter[Założenia metodologiczne i organizacja badań]{Założenia metodologiczne \\i organizacja badań}\label{ch:organizacja}

Celem niniejszego rozdziału jest przedstawienie sposobu organizacji badań porównawczych klastra Valkey i rozwiązań opartych o Redis. Rozdział opisuje przyjęte warianty badawcze, środowisko testowe, narzędzia pomiarowe, parametry eksperymentów, metryki oraz sposób analizy wyników. Szczególny nacisk położono na powtarzalność procedury, ponieważ porównanie systemów in-memory wymaga kontroli wielu czynników zewnętrznych: konfiguracji zasobów CPU i pamięci, topologii klastra, sposobu generowania ruchu, stanu danych oraz wpływu środowiska Kubernetes.

Metodyka została dobrana tak, aby umożliwić weryfikację hipotez H1--H4 sformułowanych w rozdziale~\ref{ch:wstep}. Podobny kierunek porównywania systemów in-memory, obejmujący wydajność i właściwości operacyjne, pojawia się w analizie chmurowo-natywnych magazynów danych przedstawionej przez Awaysheha i Awada~\cite{awaysheh2025}.

\section{Organizacja badań}

Badania zorganizowano jako zestaw niezależnych scenariuszy uruchamianych w Google Cloud Platform. Każdy test rozpoczyna się od przygotowania osobnego klastra Kubernetes w usłudze Google Kubernetes Engine lub od przygotowania osobnej instancji usługi zarządzanej, jeżeli dany scenariusz dotyczy Memorystore. Po zakończeniu pomiaru wyniki, logi i parametry uruchomienia są zapisywane w katalogu wynikowym, a środowisko może zostać usunięte albo odtworzone od początku dla kolejnego przebiegu.

Takie podejście ogranicza wpływ kolejności wykonywania testów na wyniki. Scenariusze operacyjne, takie jak resharding, rolling upgrade albo backup i restore, zmieniają stan klastra, tworzą lub usuwają zasoby trwałe, a w części przypadków celowo powodują niedostępność podów. Uruchamianie kolejnych testów na świeżo przygotowanym środowisku zmniejsza ryzyko, że pozostałości po poprzednim przebiegu, stare wolumeny, zmieniony stan danych albo niepełne odtworzenie klastra wpłyną na następny pomiar.

Konfiguracja klastra GKE jest utrzymywana możliwie stała w obrębie porównywanych scenariuszy. Dotyczy to przede wszystkim regionu, typu maszyn węzłów oraz sposobu uruchamiania generatorów obciążenia. Parametry klastra mogą być jednak dostosowywane do charakteru konkretnego testu. Przykładowo test backupu i odtwarzania wymaga większych wolumenów trwałych oraz innego limitu pamięci niż krótki test reshardingu. Zmiany takie są traktowane jako część definicji scenariusza i są zapisywane razem z wynikami.

Przyjęto następujący schemat pojedynczego scenariusza badawczego:

\begin{enumerate}
    \item utworzenie klastra GKE lub przygotowanie instancji Memorystore zgodnie z parametrami testu;
    \item wdrożenie klastra Valkey lub Redis 7.2 w Kubernetes albo odczyt konfiguracji usługi zarządzanej;
    \item walidacja stanu początkowego, w tym połączeń, replikacji, pokrycia slotów i gotowości klastra;
    \item przygotowanie danych testowych oraz uruchomienie generatora obciążenia, jeżeli wymaga tego scenariusz;
    \item wykonanie badanej operacji, na przykład aktualizacji, reshardingu albo odtwarzania z kopii;
    \item zapis metryk, logów, czasów faz operacji i raportów weryfikacyjnych;
    \item usunięcie albo odtworzenie środowiska przed kolejnym powtórzeniem.
\end{enumerate}

Dla każdego przebiegu zapisywane są surowe wyniki, logi oraz parametry konfiguracji (wersje oprogramowania, limity CPU i pamięci, parametry generatora obciążenia, stan klastra przed testem), co zapewnia powtarzalność badań.

\section{Warianty badawcze}

Głównym wariantem badawczym jest \textbf{Valkey Cluster w Kubernetes} --- system self-hosted oceniany pod względem wydajności, zachowania podczas operacji administracyjnych, kosztów oraz złożoności utrzymania. Punktem odniesienia wydajnościowego jest \textbf{Redis 7.2 Cluster w Kubernetes} --- klaster Redis 7.2.5 z obrazu Bitnami, wdrożony self-hosted z identyczną topologią i zasobami co Valkey, jako ostatnia wersja Redis na licencji BSD-3. Trzecim wariantem jest \textbf{Google Cloud Memorystore for Redis Cluster}, wykorzystywany jako punkt odniesienia dla kosztów infrastruktury, nakładu pracy operatora oraz operacji udostępnianych przez dostawcę chmurowego.

W przypadku wariantów self-hosted operator zarządza klastrem Kubernetes, konfiguracją Helm oraz zasobami trwałymi. W przypadku Memorystore użytkownik inicjuje operacje przez API lub narzędzia \texttt{gcloud}, natomiast wewnętrzna realizacja pozostaje ukryta --- wariant ten jest więc traktowany jako porównanie typu black-box. Tam, gdzie jest to możliwe, zachowywane są te same parametry obciążenia i podobny rozmiar danych, jednak nie zakłada się pełnej równoważności infrastruktury. Testy wykorzystujące Chaos Mesh (failover, partycja sieciowa, resilience) mogą być przeprowadzone wyłącznie dla wariantów self-hosted. Tabela~\ref{tab:testy_warianty} w dalszej części rozdziału zestawia zakres testów dla każdego wariantu.

\section{Środowisko testowe}

Ze względu na ograniczony dostęp do zasobów Google Cloud Platform (limity kwot oraz dostępnych tokenów) badania przeprowadzono w dwóch środowiskach. Testy wydajnościowe wykonano w środowisku lokalnym, natomiast testy operacyjne i odpornościowe --- w środowisku chmurowym GCP.

\subsection{Środowisko chmurowe}

Testy operacyjne i odpornościowe prowadzone są w klastrze Kubernetes utworzonym w Google Kubernetes Engine. Dla każdego scenariusza klaster GKE jest tworzony od początku, aby przebieg testu nie zależał od stanu pozostawionego przez poprzednie eksperymenty. Ze względu na ograniczenia budżetowe (dostępne kredyty GCP) stosowano dwie konfiguracje klastra GKE w zależności od tego, czy scenariusz obejmował porównanie z Memorystore.

Testy porównujące wyłącznie warianty self-hosted (failover, spójność danych, resilience) uruchamiano na mniejszym klastrze, którego konfigurację przedstawiono w listingu~\ref{lst:gke-cluster-selfhosted}.

\begin{lstlisting}[basicstyle=\ttfamily\footnotesize,columns=fullflexible,keepspaces=true,breaklines=true,breakatwhitespace=true,frame=single,framerule=0.2pt,caption={Klaster GKE dla testów self-hosted},label={lst:gke-cluster-selfhosted}]
gcloud container clusters create valkey-bench \
  --location europe-central2 \
  --node-locations europe-central2-a \
  --num-nodes 3 \
  --machine-type n2-standard-4 \
  --disk-type pd-balanced \
  --disk-size 50
\end{lstlisting}

Testy obejmujące porównanie z Memorystore (resharding, backup i restore) wymagały większych zasobów na pod, aby wyrównać parametry do usługi zarządzanej. Klaster GKE dla tych scenariuszy tworzono według listingu~\ref{lst:gke-cluster-ms}.

\begin{lstlisting}[basicstyle=\ttfamily\footnotesize,columns=fullflexible,keepspaces=true,breaklines=true,breakatwhitespace=true,frame=single,framerule=0.2pt,caption={Klaster GKE dla testów z Memorystore},label={lst:gke-cluster-ms}]
gcloud container clusters create valkey-bench \
  --location europe-central2 \
  --node-locations europe-central2-a \
  --num-nodes 2 \
  --machine-type n2-standard-8 \
  --disk-type pd-balanced \
  --disk-size 50
\end{lstlisting}

Wariant zarządzany uruchamiany był jako klaster Google Cloud Memorystore for Redis Cluster w tym samym regionie. Konfigurację usługi tworzono poleceniem~\ref{lst:memorystore-cluster-create}.

\begin{lstlisting}[basicstyle=\ttfamily\footnotesize,columns=fullflexible,keepspaces=true,breaklines=true,breakatwhitespace=true,frame=single,framerule=0.2pt,caption={Utworzenie klastra Memorystore},label={lst:memorystore-cluster-create}]
gcloud redis clusters create redis-ms \
  --region=europe-central2 \
  --network=projects/redis-vs-valkey/global/networks/default \
  --replica-count=1 \
  --node-type=redis-standard-small \
  --shard-count=3
\end{lstlisting}

\subsubsection{Normalizacja zasobów do Memorystore}\label{subsec:normalizacja}

W testach porównujących klastry self-hosted z Memorystore konieczne jest dopasowanie zasobów przydzielonych podom Valkey i Redis 7.2 do parametrów usługi zarządzanej. Typ węzła \texttt{redis-standard-small} zgodnie z dokumentacją~\cite{gcpMemorystoreNodeTypes} oferuje 2~vCPU i 6{,}5~GB pojemności całkowitej na węzeł, z czego około 5{,}2~GB stanowi domyślną pojemność użytkową (keyspace) po odjęciu zarezerwowanego narzutu systemowego.

Aby zapewnić porównywalność, klastry self-hosted w testach z udziałem Memorystore były uruchamiane z limitami i requestami zasobów na 2~vCPU i 7~GiB pamięci RAM na pod oraz z parametrem \texttt{maxmemory} ustawionym na 5200~MB. Zestawienie kluczowych parametrów dla wszystkich trzech wariantów przedstawiono w tabeli~\ref{tab:normalizacja_zasobow}.

\begin{table}[H]
\centering
\caption{Parametry zasobów w testach porównujących z Memorystore}
\label{tab:normalizacja_zasobow}
\begin{tabularx}{\textwidth}{p{3.8cm}XXX}
\toprule
\textbf{Parametr} & \textbf{Valkey} & \textbf{Redis 7.2} & \textbf{Memorystore} \\
\midrule
Obraz / typ węzła & Valkey 9.1.0 & Redis 7.2.5 (Bitnami) & redis-standard-small \\
Shardy $\times$ repliki & 3 $\times$ 1 & 3 $\times$ 1 & 3 $\times$ 1 \\
vCPU na shard & 2 & 2 & 2 \\
Pamięć na shard & 7 GiB & 7 GiB & 6{,}5 GB \\
\texttt{maxmemory} & 5200 MB & 5200 MB & $\approx$5{,}2 GB \\
Polityka eksmisji & allkeys-lru & allkeys-lru & allkeys-lru \\
\bottomrule
\end{tabularx}
\end{table}

\subsubsection{Wdrożenie klastrów self-hosted}

Klaster Valkey instalowano z wykorzystaniem chartu \texttt{valkey-helm} w wersji rozwijanej w ramach pull requestu dodającego obsługę trybu klastrowego z shardingiem~\cite{valkeyHelmPr116}. Przygotowanie repozytorium chartu przedstawiono w listingu~\ref{lst:valkey-chart-pr}.

\begin{lstlisting}[basicstyle=\ttfamily\footnotesize,columns=fullflexible,keepspaces=true,breaklines=true,breakatwhitespace=true,frame=single,framerule=0.2pt,caption={Przygotowanie chartu Valkey},label={lst:valkey-chart-pr}]
git clone https://github.com/valkey-io/valkey-helm.git
cd valkey-helm
git fetch origin pull/116/head:pr-116
git checkout pr-116
\end{lstlisting}

Instalację klastra Valkey wykonywano poleceniem \texttt{helm install} z plikiem konfiguracyjnym zawierającym topologię (3~shardy, 1~replika), wersję obrazu (9.1.0), parametry \texttt{maxmemory} oraz limity zasobów odpowiednie dla danego scenariusza. Klaster Redis 7.2 instalowano z chartu Bitnami \texttt{redis-cluster} z obrazem \texttt{bitnamilegacy/redis-cluster:7.2.5-debian-12-r4} --- ostatnią wersją Redis na licencji BSD-3-Clause. Chart Bitnami konfiguruje klaster przez parametry \texttt{cluster.nodes} i \texttt{cluster.replicas}, co daje topologię analogiczną do Valkey: 3~shardy z jedną repliką na shard, czyli 6~podów. Oba warianty self-hosted wdrażano w osobnych namespace'ach.

\subsubsection{Monitoring i chaos engineering}

Monitoring środowiska Kubernetes instalowano za pomocą chartu \texttt{kube-prometheus-stack}, który uruchamia Prometheusa i Grafanę. Testy odporności wykorzystywały Chaos Mesh do uruchamiania kontrolowanych zakłóceń w klastrze Kubernetes (runtime \texttt{containerd}).

\subsection{Środowisko lokalne}\label{subsec:srodowisko_lokalne}

Testy wydajnościowe przeprowadzono w środowisku lokalnym z wykorzystaniem klastra Kubernetes uruchomionego za pomocą narzędzia minikube. Klaster konfigurowano z trzema węzłami, aby odwzorować topologię używaną w środowisku chmurowym. Wdrożenie klastrów Valkey i Redis 7.2 odbywało się za pomocą tych samych chartów Helm i plików konfiguracyjnych co w środowisku GKE, z dostosowaniem limitów zasobów do możliwości maszyny lokalnej.

Specyfikację sprzętową maszyny lokalnej przedstawiono w tabeli~\ref{tab:spec_lokalna}, a wersje narzędzi używanych do orkiestracji testów --- w tabeli~\ref{tab:narzedzia_lokalne}.

\begin{table}[H]
\centering
\caption{Specyfikacja sprzętowa maszyny lokalnej}
\label{tab:spec_lokalna}
\begin{tabularx}{\textwidth}{p{4.0cm}X}
\toprule
\textbf{Komponent} & \textbf{Wartość} \\
\midrule
Procesor & AMD Ryzen 7 7700X, 8 rdzeni / 16 wątków, do 5{,}6 GHz \\
Pamięć RAM & 32 GB DDR5 \\
Dysk & 2$\times$ Lexar NM620 1 TB (NVMe) \\
System operacyjny & Arch Linux, jądro 7.0.9 \\
Architektura & x86\_64 \\
\bottomrule
\end{tabularx}
\end{table}

\begin{table}[H]
\centering
\caption{Wersje narzędzi w środowisku lokalnym}
\label{tab:narzedzia_lokalne}
\begin{tabularx}{\textwidth}{p{4.0cm}X}
\toprule
\textbf{Narzędzie} & \textbf{Wersja} \\
\midrule
minikube & 1.38.1 \\
kubectl & 1.35.4 \\
Helm & 4.1.4 \\
Docker & 29.2.1 \\
\bottomrule
\end{tabularx}
\end{table}

Ze względu na ograniczenia środowiska lokalnego wariant 4~vCPU na pod nie został zrealizowany. Klaster minikube z trzema węzłami dzieli zasoby jednego procesora 8-rdzeniowego, a klaster Valkey/Redis składa się z 6~podów (3~mastery + 3~repliki) plus pod generatora obciążenia. Przydzielenie 4~vCPU każdemu podowi wymagałoby łącznie co najmniej 28~vCPU wyłącznie dla podów klastra i klienta, przekraczając dostępne 16~wątków sprzętowych. Z tego powodu testy wydajnościowe ograniczono do konfiguracji 1 i 2~vCPU, co daje macierz 2~wartości CPU $\times$ 3~rozmiary payloadu $\times$ 3~proporcje operacji = 18~konfiguracji na wariant.

\section{Narzędzia badawcze}

Do generowania obciążenia wybrano narzędzie \texttt{memtier\_benchmark}. Umożliwia ono wykonywanie testów dla protokołu Redis, pracę w trybie klastrowym, sterowanie liczbą wątków i klientów, ustawienie pipeline oraz zapis wyników w formacie JSON. Jest to istotne, ponieważ wyniki muszą być możliwe do automatycznej agregacji i porównania między powtórzeniami.

W wybranych scenariuszach wykorzystywany jest również klient kontrolny uruchamiany jako osobny pod. Nie generuje on głównego obciążenia, lecz wykonuje cykliczne operacje testowe i zapisuje ich status oraz czas odpowiedzi. Dzięki temu można oddzielić przepustowość raportowaną przez \texttt{memtier\_benchmark} od obserwacji dostępności klastra z perspektywy prostego klienta aplikacyjnego.

Monitoring realizowany jest przez Prometheus, który zbiera metryki kontenerów i podów. Dane te są wykorzystywane do obliczania średniego zużycia CPU i pamięci podczas okna testowego. Wykresy oraz pliki CSV generowane są po zakończeniu testów, aby oddzielić fazę pomiaru od fazy analizy. Dzięki temu surowe wyniki mogą zostać ponownie przetworzone bez konieczności powtarzania kosztownych eksperymentów.

Do testów odpornościowych przewidziano mechanizmy chaos engineering. We wszystkich scenariuszach operacyjnych i odpornościowych zakłócenie lub operacja administracyjna jest wprowadzana po 30~s stabilizacji obciążenia, co pozwala odróżnić naturalny rozruch testu od wpływu badanego zdarzenia.

\section{Testy wydajnościowe}

Testy wydajnościowe służą weryfikacji zachowania klastra Valkey w przyjętym środowisku GKE oraz określeniu poziomu obciążenia używanego później w scenariuszach operacyjnych. Mierzone są liczba operacji na sekundę i rozkład opóźnień przy różnych rozmiarach payloadu oraz proporcjach operacji. Badane są trzy typy obciążenia: tylko odczyt, tylko zapis oraz scenariusz mieszany. Uwzględnienie różnych rozmiarów payloadu pozwala sprawdzić, czy system jest ograniczony głównie przez CPU, sieć, kopiowanie pamięci czy obsługę protokołu.

Parametry testów wydajnościowych przedstawiono w tabeli~\ref{tab:parametry_wydajnosci}. Wartości bazowe dobrano tak, aby uzyskać powtarzalny wariant eksperymentalny dla klastra self-hosted oraz punkt odniesienia dla scenariuszy operacyjnych wykonywanych pod obciążeniem. Dodatkowo można wykonywać testy saturacyjne z większą liczbą klientów i wątków, jednak powinny one być oznaczone jako osobny scenariusz, aby nie mieszać ich z wariantem bazowym.

\begin{table}[H]
\centering
\caption{Parametry bazowe testów wydajnościowych}
\label{tab:parametry_wydajnosci}
\begin{tabularx}{\textwidth}{p{4.6cm}X}
\toprule
\textbf{Parametr} & \textbf{Wartość} \\
\midrule
Narzędzie & \texttt{memtier\_benchmark} \\
Tryb klienta & \texttt{--cluster-mode} \\
Liczba vCPU na pod & 1, 2 (4 pominięto --- patrz sekcja~\ref{subsec:srodowisko_lokalne}) \\
Rozmiar payloadu & 1 KB, 10 KB, 1000 KB \\
Typ operacji & read-only \texttt{0:1}, write-only \texttt{1:0}, mixed \texttt{1:1} \\
Liczba wątków klienta & 4 \\
Liczba klientów na wątek & 25 \\
Łączna liczba połączeń & 100 \\
Pipeline & 10 \\
Maksymalna liczba komend w locie & 1000 \\
Czas pojedynczego przebiegu & 300 s \\
Liczba powtórzeń & 5 \\
Zakres kluczy & 1000000 \\
Percentyle opóźnień & p50, p95, p99, p99.9 \\
\bottomrule
\end{tabularx}
\end{table}

Wybór rozmiarów payloadu odzwierciedla rozkład rozmiarów wartości obserwowany w produkcyjnych systemach cache. Yang i współautorzy, analizując 153 klastry cache (Memcached i Redis) w Twitterze, wykazali bimodalny rozkład rozmiarów wartości: zdecydowana większość kluczy mieści się w zakresie setek bajtów do kilku kilobajtów, natomiast niewielka frakcja osiąga rozmiary rzędu setek kilobajtów~\cite{yang2020}. Trzy przyjęte punkty pomiarowe --- 1~KB, 10~KB i 1000~KB --- odpowiadają typowym, średnim i dużym wartościom w tym rozkładzie.

Planowana macierz testów obejmowała 3 wartości CPU, 3 rozmiary payloadu i 3 proporcje operacji, czyli 27 konfiguracji dla wariantu self-hosted. Ze względu na ograniczenia środowiska lokalnego opisane w sekcji~\ref{subsec:srodowisko_lokalne} zrealizowano warianty 1 i 2~vCPU, co daje 18~konfiguracji i~90~przebiegów (przy pięciu powtórzeniach). Automatyczne zapisywanie wyników oraz jednoznaczne nazewnictwo plików wynikowych pozostaje istotne ze względu na liczbę przebiegów.

\section{Testy niezawodności i odporności}

Druga grupa eksperymentów dotyczy zachowania klastra w sytuacjach awaryjnych. Testy te są związane głównie z hipotezą H3, ale ich wyniki wspierają również ocenę praktycznej użyteczności badanego rozwiązania. W środowisku produkcyjnym krótkotrwały spadek przepustowości może być akceptowalny, natomiast utrata potwierdzonych zapisów, długi brak dostępności albo lawinowy wzrost opóźnień mogą wykluczać dane rozwiązanie z zastosowań krytycznych.

Testy niezawodności, resilience oraz utrzymywalności wykorzystują wspólny profil obciążenia tła opisany w tabeli~\ref{tab:profil_operacyjny}. Różnice między scenariuszami ograniczają się do czasu przebiegu (120~s dla testów odporności, 360~s dla rolling upgrade).

\begin{table}[H]
\centering
\caption{Bazowy profil obciążenia dla testów operacyjnych}
\label{tab:profil_operacyjny}
\begin{tabularx}{\textwidth}{p{4.4cm}X}
\toprule
\textbf{Parametr} & \textbf{Wartość} \\
\midrule
Protokół i tryb pracy & \texttt{redis}, \texttt{--cluster-mode} \\
Liczba wątków & 4 \\
Liczba klientów na wątek & 16 \\
Łączna liczba połączeń & 64 \\
Czas pojedynczego przebiegu & 120 s (360 s dla rolling upgrade) \\
Proporcja operacji & \texttt{1:1} (\texttt{SET}:\texttt{GET}) \\
Pipeline & 1 \\
Rozmiar payloadu & 1024 B \\
Zakres kluczy & 0--100000 \\
Percentyle opóźnień & p50, p95, p99, p99.9 \\
Liczba powtórzeń & 5 \\
\bottomrule
\end{tabularx}
\end{table}

\subsection{Failover mastera}\label{subsec:failover}

Test failoveru polega na ubiciu jednego poda pełniącego rolę mastera podczas ciągłego obciążenia i obserwacji, jak klaster promuje replikę oraz jak klient odczuwa degradację. Awaria jest wstrzykiwana za pomocą zasobu \texttt{PodChaos} w trybie \texttt{pod-failure} generowanego dynamicznie przez skrypt benchmarkowy.

Klientem obciążeniowym jest dedykowany program napisany w Pythonie, oparty na bibliotece \texttt{redis-py} w trybie klastrowym. W~odróżnieniu od \texttt{memtier\_benchmark}, klient ten kontynuuje wysyłanie operacji do zdrowych shardów nawet po utracie połączenia z jednym masterem, dzięki czemu zarejestrowana przepustowość odzwierciedla faktyczną dostępność klastra, a nie artefakt globalnego refresh topologii klienta. Klient uruchamia 64~współbieżne wątki robocze (\texttt{threads=4}, \texttt{clients=16}), każdy z własnym połączeniem do klastra, przy payloadzie 1~KB i proporcji operacji 1:1. Timeout socketu wynosi 1~s, a po błędzie połączenia wątek tworzy nowe połączenie z~backoffem 10~ms.

Procedura każdego powtórzenia przebiega następująco:

\begin{enumerate}
    \item uruchomienie poda z klientem obciążeniowym na 120~s;
    \item stabilizacja obciążenia przez 30~s;
    \item identyfikacja aktualnych masterów w klastrze poleceniem \texttt{CLUSTER NODES};
    \item wygenerowanie i zastosowanie manifestu \texttt{PodChaos} ubijającego wybranego mastera z parametrem \texttt{gracePeriod: 1} na czas 60~s;
    \item kontynuacja pomiaru do zakończenia okna 120~s;
    \item zapis raportu JSON z per-sekundowymi statystykami sukcesu, błędów i opóźnień.
\end{enumerate}

Wariant podstawowy obejmuje awarię jednego mastera z trzech, ponieważ taki scenariusz jest najłatwiejszy do powtórzenia i interpretacji. Dla każdego wariantu wykonano po pięć powtórzeń.

Mierzone metryki obejmują: przepustowość w oknie bazowym (30~s przed chaos), przepustowość w oknie chaos (60~s trwania \texttt{PodChaos}), procentowy spadek przepustowości, czas okna błędów klienta (od pierwszego do ostatniego błędu połączenia), szacunkową liczbę utraconych operacji, szczytowe opóźnienie p99 podczas failoveru, liczbę błędów połączenia i odpowiedzi \texttt{CLUSTERDOWN}.

Kluczową metryką dla hipotezy H3 jest \textit{czas okna błędów} (\textit{error window}), definiowany jako liczba sekund od pierwszego do ostatniego błędu połączenia zarejestrowanego przez klienta. Metryka ta izoluje rzeczywisty czas od utraty dostępności mastera do momentu, w którym klient uzyskuje połączenie z wypromowaną repliką. Jest niezależna od czasu trwania zasobu \texttt{PodChaos} (60~s) oraz od ewentualnego opóźnienia terminacji poda w Kubernetes, ponieważ obejmuje wyłącznie okres, w którym klient obserwuje błędy po stronie dotkniętego sharda.

\subsection{Spójność danych podczas partycji sieciowej}\label{subsec:split_brain}

Test spójności danych bada zachowanie klastra podczas partycji sieciowej (ang.~\textit{split-brain}) wprowadzonej za pomocą zasobu \texttt{NetworkChaos} z Chaos Mesh.

Procedura każdego powtórzenia przebiega następująco:

\begin{enumerate}
    \item skrypt odczytuje topologię klastra poleceniem \texttt{CLUSTER NODES} i wybiera mastera jednego sharda jako stronę mniejszościową (minority); jego replika pozostaje po stronie większościowej (majority);
    \item pody klastra są etykietowane: wybrany master otrzymuje etykietę \texttt{chaos-side=minority}, pozostałe pody (w tym replika tego mastera) --- etykietę \texttt{chaos-side=majority};
    \item uruchamiany jest klient konsystencji, który zapisuje klucze w formacie \texttt{sb.RUN.cID.SEQ} bez hash tagów, dzięki czemu klucze rozkładają się równomiernie po wszystkich slotach; klient rejestruje, które klucze trafiają do slotów mastera mniejszościowego, a które do masterów większościowych;
    \item po stabilizacji skrypt aplikuje manifest \texttt{NetworkChaos} realizujący obustronną partycję sieciową między podami oznaczonymi \texttt{chaos-side=minority} a podami \texttt{chaos-side=majority}; partycja trwa 90~s;
    \item w oknie partycji master mniejszościowy może przez czas nieprzekraczający \texttt{cluster-node-timeout} nadal przyjmować i potwierdzać zapisy; po upływie timeoutu zwraca \texttt{CLUSTERDOWN}, a strona większościowa promuje replikę na nowego mastera;
    \item po zakończeniu partycji (heal) dawny master mniejszościowy wraca jako replika i synchronizuje dane z nowym masterem; klucze potwierdzone wyłącznie przez dawnego mastera mniejszościowego w oknie partycji mogą zostać utracone;
    \item po 15~s stabilizacji klient weryfikuje wszystkie potwierdzone klucze i generuje raport JSON.
\end{enumerate}

Mierzone metryki obejmują: liczbę utraconych potwierdzonych kluczy z podziałem na klucze ze slotów mniejszościowych (\texttt{keys\_missing\_minority}) i większościowych (\texttt{keys\_missing\_majority}), wskaźnik utraty danych (\texttt{loss\_rate}), wskaźnik utraty na stronie mniejszościowej (\texttt{minority\_loss\_rate}), liczbę nieudanych żądań zapisu oraz odsetek żądań dotkniętych awarią. Oczekiwany wynik to utrata kluczy wyłącznie ze strony mniejszościowej, co potwierdza znane właściwości modelu spójności Redis/Valkey Cluster.

\subsection{Odporność na niedobór zasobów (resilience)}\label{subsec:resilience}

Testy odporności (ang.~\textit{resilience}) badają zachowanie klastra podczas niedoboru zasobów systemowych. Zbadano dwa scenariusze: obciążenie CPU masterów oraz presję pamięci z polityką eksmisji. Oba scenariusze są ograniczone do wariantów self-hosted, ponieważ wymagają wstrzykiwania awarii na poziomie podów Kubernetes. Dla każdej konfiguracji wykonano po pięć powtórzeń.

\subsubsection{Obciążenie CPU}

Scenariusz obciążenia CPU polega na wstrzyknięciu stresu procesora na wybrane pody pełniące rolę mastera za pomocą zasobu \texttt{StressChaos} z Chaos Mesh. Procedura każdego powtórzenia przebiega następująco:

\begin{enumerate}
    \item uruchomienie poda z \texttt{memtier\_benchmark} zgodnie z profilem z tabeli~\ref{tab:profil_operacyjny};
    \item stabilizacja obciążenia przez 30~s;
    \item identyfikacja aktualnych masterów w klastrze poleceniem \texttt{CLUSTER NODES} i wybranie zadanej liczby podów docelowych;
    \item wygenerowanie i zastosowanie manifestu \texttt{StressChaos} uruchamiającego 4~procesy obciążające CPU (\texttt{workers: 4}, \texttt{load: 100}) na wybranych masterach przez 30~s;
    \item kontynuacja pomiaru przez \texttt{memtier\_benchmark} do zakończenia okna 120~s;
    \item zapis wyników JSON, logów oraz pliku z parametrami obciążenia.
\end{enumerate}

Test wykonano w trzech wariantach: obciążenie jednego, dwóch oraz trzech masterów (z trzech dostępnych w klastrze). Pozwala to zaobserwować, jak degradacja rośnie w zależności od odsetka dotkniętych shardów.

Mierzone metryki obejmują: spadek przepustowości względem poziomu bazowego (próg degradacji 80\%), czas trwania degradacji, minimalne ops/sec w oknie obciążenia, szczytowe opóźnienie p99 i p99.9 podczas stresu, bazowy poziom ops/sec i p99 przed obciążeniem oraz informację, czy po zakończeniu obciążenia nastąpił powrót do co najmniej 90\% poziomu bazowego.

\subsubsection{Presja pamięci}

Scenariusz presji pamięci bada zachowanie klastra po wypełnieniu pamięci do limitu \texttt{maxmemory} z aktywną polityką eksmisji \texttt{allkeys-lru}. Procedura każdego powtórzenia przebiega następująco:

\begin{enumerate}
    \item konfiguracja polityki eksmisji na wszystkich węzłach klastra poleceniem \texttt{CONFIG SET maxmemory-policy allkeys-lru};
    \item zapisanie stanu klastra przed wypełnieniem (zużycie pamięci, liczba kluczy, liczba eksmisji);
    \item wypełnienie klastra za pomocą dedykowanego skryptu seedującego liczbą kluczy wystarczającą do osiągnięcia 100\% limitu \texttt{maxmemory} (3\,145\,728 kluczy o rozmiarze 1~KB przy łącznym limicie 3072~MB dla całego klastra, czyli po 1024~MB na shard);
    \item zapisanie stanu klastra po wypełnieniu (zużycie pamięci, liczba kluczy, liczba eksmisji, błędy OOM);
    \item uruchomienie poda z \texttt{memtier\_benchmark} na 30~s z obciążeniem mieszanym zgodnym z profilem z tabeli~\ref{tab:profil_operacyjny};
    \item zapisanie stanu klastra po zakończeniu obserwacji;
    \item zapis raportu JSON zawierającego wszystkie pomiary, raport seedowania oraz wyniki \texttt{memtier\_benchmark}.
\end{enumerate}

Między kolejnymi powtórzeniami klaster jest czyszczony poleceniem \texttt{FLUSHALL}, aby zapewnić powtarzalny punkt startowy. Limit \texttt{maxmemory} wynosi 1024~MB na shard, czyli 3072~MB łącznie dla klastra trzech shardów. Wszystkie raportowane wielkości (liczba kluczy, eksmisje, wykorzystana pamięć) podawane są sumarycznie dla całego klastra.

Mierzone metryki obejmują: liczbę eksmitowanych kluczy z podziałem na fazę wypełniania i fazę obserwacji, liczbę kluczy pozostających po teście (\texttt{dbsize}), wykorzystaną pamięć po teście, przepustowość i opóźnienie p99 w oknie obserwacji, liczbę błędów OOM, liczbę błędów zapisu, czas seedowania oraz informację, czy osiągnięto 100\% limitu \texttt{maxmemory}.

\section{Testy utrzymywalności}\label{sec:testy_utrzymywalnosci}

Testy utrzymywalności służą weryfikacji hipotezy H4. Ich celem jest ocena, ile pracy operacyjnej wymaga utrzymanie klastra self-hosted w porównaniu z usługą zarządzaną. W tej części obok metryk technicznych mierzone są również liczba kroków, liczba komend, czas wykonania procedury, konieczność ręcznej interwencji oraz ryzyko błędnej konfiguracji. W rozdziale z wynikami przedstawiono dwa scenariusze wykonywane pod obciążeniem: rolling upgrade oraz resharding klastra.

Obciążenie tła w scenariuszach utrzymaniowych jest zgodne z profilem z tabeli~\ref{tab:profil_operacyjny}. Dla reshardingu czas przebiegu wynosi 120~s (5~powtórzeń na fazę i metodę), natomiast rolling upgrade wykorzystuje okno 360~s z 5~powtórzeniami.

\subsection{Rolling upgrade}\label{subsec:rolling_upgrade}

Test rolling upgrade obejmuje aktualizację wersji klastra Valkey podczas ciągłego obciążenia. Każde powtórzenie rozpoczyna się od przywrócenia wersji bazowej, dzięki czemu kolejne przebiegi startują z tego samego stanu. Następnie skrypt uruchamia pod z \texttt{memtier\_benchmark} oraz osobny pod klienta kontrolnego. Klient kontrolny działa niezależnie od generatora obciążenia i cyklicznie wykonuje sekwencję \texttt{SET}, \texttt{GET} i \texttt{DEL} w trybie cluster-aware. Próby wykonywane są co 0{,}2~s, a limit czasu odpowiedzi pojedynczej operacji wynosi 2~s.

Po stabilizacji obciążenia wykonywany jest rzeczywisty \texttt{helm upgrade}, który zmienia wersję obrazu Valkey i wymusza rolling update podów StatefulSetu. Czas operacji administracyjnej liczony jest od rozpoczęcia \texttt{helm upgrade} do zakończenia rollout oraz pozytywnej kontroli gotowości klastra z perspektywy klienta. Po zakończeniu testu zapisywane są wyniki \texttt{memtier\_benchmark}, logi oraz plik CSV klienta kontrolnego zawierający znacznik czasu, czas odpowiedzi, status próby i treść błędu.

\subsection{Skalowanie horyzontalne i resharding}\label{subsec:resharding}

Test reshardingu obejmuje zmianę liczby shardów z trzech do czterech oraz powrót z czterech do trzech shardów. Dla klastra Valkey porównywane są dwa warianty migracji slotów: klasyczny tryb \texttt{valkey-cli --cluster rebalance} oraz tryb atomowy uruchamiany flagą \texttt{--cluster-use-atomic-slot-migration}. Dla klastra Redis 7.2 wykonywana jest klasyczna migracja slotów w tej samej topologii Kubernetes. Mechanizm atomowej migracji slotów jest dostępny w Valkey od wersji 9.0~\cite{valkey90release}; w badaniu wykorzystano obraz Valkey 9.1~\cite{valkey91release}, ponieważ używane narzędzie \texttt{valkey-cli} udostępnia w tej wersji odpowiednią flagę. Na potrzeby benchmarku chart Helm został rozszerzony o możliwość wyłączenia automatycznego równoważenia slotów oraz przekazania flagi atomowej migracji slotów do operacji reshardingu~\cite{valkeyHelmPr116}.

Przed rozpoczęciem pomiaru skrypt sprawdza, czy klaster ma trzy mastery, pełne pokrycie 16384 slotów i stan \texttt{cluster\_state:ok}. Następnie uruchamiany jest pod z \texttt{memtier\_benchmark}; po stabilizacji obciążenia rozpoczyna się operacja. Dla skalowania w górę w wariantach self-hosted wykonywany jest \texttt{helm upgrade} zwiększający liczbę shardów do czterech. Automatyczne równoważenie chartu jest wyłączone, aby migracja slotów była wykonywana jawnie przez skrypt benchmarkowy. Istniejące pody są zabezpieczane przed rolling restartem przez ustawienie partycji aktualizacji StatefulSetu.

Po utworzeniu nowych podów skrypt czeka, aż będą gotowe oraz widoczne w metadanych klastra. Następnie wykonywane jest jawne przeniesienie slotów na dodatkowy master. Dla Valkey legacy używana jest klasyczna migracja slotów, dla Redis 7.2 klasyczna migracja wykonywana narzędziem \texttt{redis-cli}, a dla Valkey atomic ta sama operacja uruchamiana jest z flagą atomowej migracji. Po zakończeniu migracji skrypt ponownie sprawdza stan klastra i zapisuje czasy faz: tworzenie nowych podów, migrację slotów oraz czas całej operacji.

Faza skalowania w dół rozpoczyna się osobnym przebiegiem \texttt{memtier\_benchmark}. Po 30~s stanu ustalonego skrypt przenosi sloty z dodatkowego sharda na pozostałe mastery, usuwa dodatkowe węzły z metadanych klastra poleceniem \texttt{del-node}, a następnie wykonuje \texttt{helm upgrade} przywracający \texttt{cluster.shards=3}. Mierzony czas obejmuje więc zarówno migrację slotów, jak i część administracyjną potrzebną do usunięcia sharda oraz przywrócenia stabilnej konfiguracji.

Dla Memorystore procedura ma charakter black-box. Skrypt uruchamia pod z \texttt{memtier\_benchmark}, po stabilizacji wywołuje zmianę liczby shardów przez \texttt{gcloud redis clusters update --shard-count}. Pomiar kończy się dopiero po tym, gdy opis klastra zwraca oczekiwaną liczbę shardów i stan gotowości usługi. Wariant zarządzany nie pozwala rozdzielić czasu tworzenia zasobów, wewnętrznej migracji slotów i kontroli stanu, dlatego w wynikach jest prezentowany jako pojedynczy blok czasu.

\subsection{Backup i restore}\label{subsec:backup_restore}

Test backupu i odtwarzania obejmuje wypełnienie klastra danymi, utworzenie kopii zapasowej oraz pełne odtworzenie klastra z tej kopii. Procedura jest wykonywana pięciokrotnie dla każdego wariantu, aby uzyskać powtarzalne wyniki. Mierzone są czas poszczególnych faz operacji, całkowity czas odtworzenia klastra oraz zgodność liczby kluczy po odtworzeniu.

Dla klastrów Valkey i Redis 7.2 w Kubernetes przyjęto metodę online opartą o repliki oraz snapshoty wolumenów GCE Persistent Disk. Każde powtórzenie rozpoczyna się od wypełnienia klastra danymi testowymi z losowymi wartościami. Następnie skrypt wybiera po jednej replice dla każdego mastera, wykonuje na tych replikach \texttt{BGSAVE}, waliduje utworzone pliki RDB i tworzy snapshoty PVC replik. Klaster źródłowy nie jest skalowany do zera podczas backupu, dzięki czemu pomiar obejmuje procedurę, która może być wykonywana przy działającym klastrze.

Faza odtwarzania polega na utworzeniu świeżego klastra docelowego, przywróceniu dysków ze snapshotów replik, zamontowaniu ich jako tymczasowych PVC oraz skopiowaniu plików RDB wewnątrz Kubernetes do PVC masterów klastra docelowego. PVC replik są czyszczone przed startem, aby repliki odtworzyły stan przez synchronizację z masterami. Skrypt czeka następnie na osiągnięcie stanu \texttt{cluster\_state:ok} z co najmniej trzema masterami. Mierzone czasy obejmują między innymi: czas \texttt{BGSAVE}, walidację RDB, tworzenie snapshotów, tworzenie świeżego klastra, tworzenie dysków ze snapshotów, kopiowanie RDB, start podów oraz czas osiągnięcia gotowości klastra.

Dla usługi Memorystore procedura wykorzystuje zarządzany mechanizm backupu dostawcy chmurowego. Po wypełnieniu klastra danymi testowymi skrypt tworzy kopię zapasową poleceniem \texttt{gcloud redis clusters create-backup}. Następnie odtworzenie realizowane jest przez utworzenie nowego klastra z parametrem \texttt{--import-managed-backup}, wskazującym zasób backupu. Skrypt czeka na osiągnięcie stanu gotowości nowego klastra, po czym uruchamia pod weryfikacyjny łączący się z odtworzonym klastrem i sprawdzający liczbę kluczy oraz rozmiar danych. Po zakończeniu weryfikacji klaster odtworzeniowy jest usuwany. Wariant zarządzany nie pozwala rozdzielić wewnętrznych faz operacji, dlatego w wynikach prezentowany jest łączny czas od wywołania komendy tworzenia backupu do zakończenia operacji oraz łączny czas od utworzenia klastra z backupu do jego pełnej dostępności.

\section{Metryki i sposób analizy}

Zestaw metryk dobrano tak, aby umożliwić ocenę zarówno wydajności, jak i niezawodności oraz utrzymywalności. Tabela~\ref{tab:metryki_badawcze} przedstawia główne kategorie metryk.

\begin{table}[H]
\centering
\caption{Główne metryki wykorzystywane w analizie}
\label{tab:metryki_badawcze}
\begin{tabularx}{\textwidth}{p{3.8cm}X}
\toprule
\textbf{Kategoria} & \textbf{Metryki} \\
\midrule
Wydajność & Ops/sec, średnia z powtórzeń, odchylenie standardowe, stosunek przepustowości wariantów. \\
Opóźnienia & p50, p95, p99, p99.9. \\
Zasoby & Średnie zużycie CPU, zużycie CPU per pod, zużycie pamięci, restarty podów (z Prometheusa). \\
Failover & Czas okna błędów klienta (error window), utracone operacje, spadek przepustowości w oknie chaos, nieudane żądania, odpowiedzi \texttt{CLUSTERDOWN}. \\
Spójność danych & Utracone potwierdzone zapisy, wskaźnik utraty (loss rate), minority/majority loss rate, zapisy odrzucone przez klaster. \\
Resilience & Spadek przepustowości [\%], min ops/sec, szczytowe p99, eksmisje, błędy OOM, powrót do poziomu bazowego (baseline). \\
Utrzymywalność & Liczba kroków procedury, liczba komend operatora, czas procedury, odsetek błędów klienta kontrolnego (error rate). \\
Backup i restore & Czas faz (BGSAVE, snapshot PVC, kopiowanie RDB, odzyskanie klastra), zgodność liczby kluczy. \\
Koszty & Koszt miesięczny infrastruktury, koszt na 100 tys. ops/sec, próg opłacalności self-hosted vs~managed. \\
\bottomrule
\end{tabularx}
\end{table}

Dla każdej konfiguracji liczbowej wykonywanych jest co najmniej pięć powtórzeń. Następnie obliczana jest średnia arytmetyczna oraz odchylenie standardowe. W scenariuszach porównujących warianty obliczana jest również różnica procentowa między nimi --- w zależności od testu dotyczy to porównania Valkey z Redis 7.2, Valkey z Memorystore lub wszystkich trzech wariantów jednocześnie.

Aby ocenić, czy zaobserwowane różnice nie wynikają wyłącznie z losowej zmienności między powtórzeniami, w scenariuszach porównawczych przeprowadzany jest dwustronny test Manna-Whitneya~U na poziomie istotności $\alpha = 0{,}05$. Test ten nie wymaga założenia o normalności rozkładu, co jest istotne przy małej próbie ($N = 5$), dla której test Shapiro-Wilka ma bardzo niską moc~\cite{shapiro1965}. Jako uzupełnienie podawany jest test Welcha (parametryczny odpowiednik t-Studenta bez założenia równych wariancji). Jako miarę wielkości efektu raportowana jest korelacja rang-biserialna $r$ (naturalna miara efektu dla testu Manna-Whitneya; $|r| < 0{,}3$ oznacza mały efekt, $0{,}3 \leq |r| < 0{,}5$ --- średni, $|r| \geq 0{,}5$ --- duży). W scenariuszach z wieloma jednoczesnymi porównaniami (np.~18~konfiguracji wydajnościowych) stosowana jest korekta Benjaminiego-Hochberga (FDR), aby kontrolować odsetek fałszywych odkryć. Przy porównaniu trzech lub więcej grup (np.~backup Valkey vs Redis vs Memorystore) jako test omnibusowy stosowany jest test Kruskala-Wallisa, a następnie porównania parami z korektą na wielokrotne testowanie.

Ze względu na ograniczoną liczbę powtórzeń wyniki należy interpretować ostrożnie i traktować je jako eksperymentalne porównanie w określonym środowisku, a nie jako uniwersalną charakterystykę wszystkich możliwych wdrożeń. Wynik $p \geq 0{,}05$ nie oznacza, że różnica nie istnieje --- oznacza jedynie brak wystarczających dowodów na jej powtarzalność przy danej liczbie obserwacji.

Wyniki prezentowane są w formie tabel i wykresów. Dla testów wydajnościowych stosowane są wykresy słupkowe ops/sec, wykresy percentyli opóźnień oraz heatmapy zależności przepustowości (ang.~\textit{throughput}) od CPU i rozmiaru payloadu. Dla testów failover, resilience, upgrade i resharding stosowane są wykresy czasowe, na których oznaczone jest okno wprowadzenia zakłócenia lub operacji administracyjnej.

\section{Zakres testów dla poszczególnych wariantów}

Nie wszystkie testy mogą być przeprowadzone dla każdego z trzech wariantów badawczych. Tabela~\ref{tab:testy_warianty} zestawia zakres testów dla Valkey, Redis 7.2 i Memorystore. Ograniczenia wynikają przede wszystkim z braku możliwości stosowania Chaos Mesh wobec usługi zarządzanej oraz z faktu, że wybrane testy operacyjne (rolling upgrade, backup PVC) dotyczą specyficznych mechanizmów jednego wariantu.

\begin{table}[H]
\centering
\caption{Zakres testów dla poszczególnych wariantów badawczych}
\label{tab:testy_warianty}
\begin{tabular}{lccc}
\toprule
\textbf{Test} & \textbf{Valkey} & \textbf{Redis 7.2} & \textbf{Memorystore} \\
\midrule
Wydajność          & \cmark & \cmark & \xmark \\
Resharding         & \cmark & \cmark & \cmark \\
Rolling upgrade    & \cmark & \xmark & \xmark \\
Backup i restore   & \cmark & \cmark & \cmark \\
Failover mastera   & \cmark & \cmark & \xmark \\
Spójność danych    & \cmark & \cmark & \xmark \\
Resilience         & \cmark & \cmark & \xmark \\
\bottomrule
\end{tabular}
\end{table}

Testy wykorzystujące Chaos Mesh (failover mastera, spójność danych, resilience) wymagają dostępu do podów i mogą być przeprowadzone wyłącznie dla wariantów self-hosted. W przypadku Memorystore dostawca nie udostępnia bezpośredniego dostępu do węzłów, co uniemożliwia wstrzyknięcie nagłej awarii mastera porównywalnej z testem \texttt{pod-failure}. Memorystore for Redis Cluster nie udostępnia również komendy ręcznego failoveru pojedynczego sharda~\cite{gcpMemorystoreHA} --- jedyną dostępną operacją jest symulacja planowanego zdarzenia konserwacji~\cite{gcpMemorystoreSimulateMaintenance}, która wywołuje skoordynowany, bezstratny failover na wszystkich węzłach i nie jest równoważna nagłej awarii. Z tego powodu druga część hipotezy H3, dotycząca porównywalności z mechanizmami Memorystore, nie może zostać zweryfikowana eksperymentalnie w~ramach niniejszej pracy.

\section{Kryteria interpretacji wyników}

Ocena wariantu Valkey jest przeprowadzana osobno dla każdego scenariusza. W scenariuszach wydajnościowych korzystniejszy wynik oznacza wyższą przepustowość lub niższe opóźnienia bez pogorszenia stabilności i zużycia zasobów. Pojedynczy lepszy wynik ops/sec nie jest wystarczający, jeżeli towarzyszy mu istotny wzrost p99 lub p99.9 latency. W scenariuszach operacyjnych równie ważne są czas procedury, liczba wymaganych kroków oraz ryzyko błędu operatora.

W przypadku hipotezy kosztowej porównanie obejmuje nie tylko koszt infrastruktury, ale także koszt pracy operacyjnej. Klaster self-hosted może być tańszy pod względem zasobów obliczeniowych, ale wymaga samodzielnego monitorowania, aktualizacji, reagowania na awarie oraz utrzymywania procedur backup/restore. Z tego względu wariant Valkey self-hosted zostanie uznany za korzystniejszy kosztowo tylko wtedy, gdy oszczędność infrastrukturalna nie zostanie zniwelowana przez nakład pracy utrzymaniowej (ang.~\textit{maintenance overhead}) oraz ryzyko operacyjne.

Ostateczna ocena zostanie przeprowadzona osobno dla każdej hipotezy H1--H4 w rozdziale~\ref{ch:podsumowanie}. Pozwoli to oddzielić wnioski dotyczące wydajności od wniosków dotyczących kosztów i utrzymywalności.
